\chaptertitle{5}{第五章 向量代数与空间解析几何}

\begin{summarybox}{本章学习目标}
第五章主要把平面几何的“点、线、面”推广到空间中。复习时抓住两条主线：一是用向量运算判断平行、垂直、共面；二是用法向量和方向向量写出平面、直线、曲面与曲线的方程。
\end{summarybox}

\begin{methodbox}{本章常见题型}
\begin{itemize}
  \item 用数量积判断垂直，用向量积判断平行或求法向量。
  \item 写平面方程：通常先找一点和法向量。
  \item 写直线方程：通常先找一点和方向向量。
  \item 求距离：先识别点到平面、点到直线，还是异面直线距离。
  \item 识别柱面、旋转曲面和二次曲面。
\end{itemize}
\end{methodbox}

\subsection{第一节 向量及其线性关系}
 \begin{dl}{空间直角坐标系}{}
以空间某一个定点 $O$ 为原点，作三条相互垂直的数轴，分别叫做 $x$ 轴、$y$ 轴和 $z$ 轴。三轴正向符合右手法则，由此构成空间直角坐标系。三个坐标平面把空间分成八个卦限。
 \end{dl}
\begin{center}
\begin{tikzpicture}[mpdiagram,scale=0.92]
  \coordinate (O) at (0,0);
  \coordinate (X) at (3.5,-0.78);
  \coordinate (Y) at (3.0,0.86);
  \coordinate (Z) at (0,2.7);
  \coordinate (A) at (2.35,-0.52);
  \coordinate (B) at (1.95,0.56);
  \coordinate (Q) at (4.30,0.04);
  \coordinate (P) at (4.30,2.04);

  \draw[mpaxis] (O) -- (X) node[below right,mplabel] {$x$};
  \draw[mpaxis] (O) -- (Y) node[right,mplabel] {$y$};
  \draw[mpaxis] (O) -- (Z) node[above,mplabel] {$z$};
  \draw[mpguide] (A) -- (Q);
  \draw[mpguide] (B) -- (Q);
  \draw[mpguide] (Q) -- (P);
  \draw[mpguide] (A) -- ++(0,2.0);
  \draw[mpguide] (B) -- ++(2.35,1.48);
  \fill[MaterialRed] (P) circle (1.25pt) node[above right,mplabel] {$P(x,y,z)$};
  \fill[MaterialBlueGrey] (Q) circle (1pt) node[below right,mplabel] {$(x,y,0)$};
  \node[below left,mplabel] at (O) {$O$};
\end{tikzpicture}
\diagramcaption{空间点的三个坐标可以看成先落到 $xOy$ 平面，再沿 $z$ 方向升高。}
\end{center}
	  \begin{dl}{向量}{}
向量又称矢量，是既有大小、又有方向的量，通常记作 $\vec{a}$、$\boldsymbol{a}$ 或 $\overrightarrow{M_1 M_2}$。
    若向量
    $\boldsymbol{a}=(x,y,z)=x\boldsymbol{i}+y\boldsymbol{j}+z\boldsymbol{k}$,则其模$|\boldsymbol{a}|=\sqrt{x^2+y^2+z^2}$,与$\boldsymbol{a}$同方向的单位向量\[\boldsymbol{e}_a=\frac{\boldsymbol{a}}{|\boldsymbol{a}|}=(\cos\alpha,\cos\beta,\cos\gamma),\]其中$\cos\alpha,\cos\beta,\cos\gamma$称为向量$\boldsymbol{a}$的方向余弦
 \end{dl}
   \begin{zs}{以$A(x_1,y_1,z_1)$为起点，$B(x_2,y_2,z_2)$为终点的向量}
 \[\overrightarrow{AB}=(x_2-x_1,y_2-y_1,z_2-z_1),\]
 空间两点$A(x_1,y_1,z_1)$与$B(x_2,y_2,z_2)$间的直线距离为
 \[|\overrightarrow{AB}|=\sqrt{(x_2-x_1)^2+(y_2-y_1)^2+(z_2-z_1)^2}.\]
   \end{zs}
 \begin{zs}{向量$\boldsymbol{b}$在向量$\boldsymbol{a}$上的投影$\text{Prj}_{\boldsymbol{a}}\boldsymbol{b}=|\boldsymbol{b}|\cos\theta$,投影向量为 $|\boldsymbol{b}|\cos\theta \cdot \boldsymbol{e}_a$, 其中 $\theta$ 为两向量之间的夹角, $\boldsymbol{e}_a$ 为与 $\boldsymbol{a}$ 同方向的单位向量.}
   \end{zs}
   \begin{zs}
 {自由向量是具有大小和方向、不考虑起点位置的向量; 单位向量是模等于1的向量; 零向量是模等于零的向量, 一般记作 $\vec{0}$ 或者 $\boldsymbol{0}$, 零向量的方向是任意的; $\boldsymbol{a}$ 的负向量是与 $\boldsymbol{a}$ 大小相等、方向相反的向量, 记作 $-\boldsymbol{a}$; 两向量相等是指大小相等、方向相同; 向量平行, 也称向量共线, 是指两个非零向量 $\boldsymbol{a}$ 和 $\boldsymbol{b}$ 方向相同或相反, 记作 $\boldsymbol{a} \parallel \boldsymbol{b}$, 零向量与任何向量都平行.}
        
    \end{zs}
    \begin{dl}{向量的线性运算}{}
设向量 $\boldsymbol{a} = (x_1,y_1,z_1)$, $\boldsymbol{b} = (x_2,y_2,z_2)$, 则
\[
\boldsymbol{a} \pm \boldsymbol{b} = (x_1 \pm x_2,\ y_1 \pm y_2,\ z_1 \pm z_2), \quad
\lambda\boldsymbol{a} = (\lambda x_1,\ \lambda y_1,\ \lambda z_1),
\]
其中 $\lambda$ 为任意实数.
\end{dl}
\begin{zs}
结论: 设向量 $\boldsymbol{a} \neq \boldsymbol{0}$, 则非零向量 $\boldsymbol{b}$ 平行于 $\boldsymbol{a} \Leftrightarrow$ 存在实数 $\lambda$, 使得 $\boldsymbol{b} = \lambda\boldsymbol{a}$ (亦称向量 $\boldsymbol{a}$ 与 $\boldsymbol{b}$ 共线).
\end{zs}
\begin{dl}{ 数量积}{}
设向量 $\boldsymbol{a} = (x_1,y_1,z_1)$, $\boldsymbol{b} = (x_2,y_2,z_2)$, 则
\[
\boldsymbol{a} \cdot \boldsymbol{b} = |\boldsymbol{a}||\boldsymbol{b}|\cos(\widehat{\boldsymbol{a},\boldsymbol{b}}) = x_1x_2 + y_1y_2 + z_1z_2.
\]
\end{dl}
 \begin{zs}
运算律:
\[
\boldsymbol{a} \cdot \boldsymbol{b} = \boldsymbol{b} \cdot \boldsymbol{a},\quad
\boldsymbol{a} \cdot (\boldsymbol{b} + \boldsymbol{c}) = \boldsymbol{a} \cdot \boldsymbol{b} + \boldsymbol{a} \cdot \boldsymbol{c},\quad
\lambda(\boldsymbol{a} \cdot \boldsymbol{b}) = (\lambda\boldsymbol{a}) \cdot \boldsymbol{b} = \boldsymbol{a} \cdot (\lambda\boldsymbol{b}).
\]
\end{zs}
 \begin{zs}
非零向量夹角余弦:
\[
\cos(\widehat{\boldsymbol{a},\boldsymbol{b}}) = \frac{\boldsymbol{a} \cdot \boldsymbol{b}}{|\boldsymbol{a}||\boldsymbol{b}|}.
\]
\end{zs}
 \begin{zs}
向量的模:
\[
|\boldsymbol{a}|^2 = \boldsymbol{a} \cdot \boldsymbol{a},\quad \text{即} \quad |\boldsymbol{a}| = \sqrt{\boldsymbol{a} \cdot \boldsymbol{a}}.
\]
\end{zs}
 \begin{zs}
向量的投影:
\[
\text{Prj}_{\boldsymbol{a}} \boldsymbol{b} = \frac{\boldsymbol{a} \cdot \boldsymbol{b}}{|\boldsymbol{a}|}.
\]
\end{zs}
 \begin{zs}
应用: 判定非零向量垂直.
\[
\boldsymbol{a} \perp \boldsymbol{b} \Leftrightarrow \boldsymbol{a} \cdot \boldsymbol{b} = 0 \Leftrightarrow x_1x_2 + y_1y_2 + z_1z_2 = 0.
\]
\end{zs}
\begin{dl}{ 向量积}{}
设向量 $\boldsymbol{a} = (x_1,y_1,z_1)$，$\boldsymbol{b} = (x_2,y_2,z_2)$，则向量积
\[
\boldsymbol{c} = \boldsymbol{a} \times \boldsymbol{b} =
\begin{vmatrix}
\boldsymbol{i} & \boldsymbol{j} & \boldsymbol{k} \\
x_1 & y_1 & z_1 \\
x_2 & y_2 & z_2
\end{vmatrix},
\]
其中 $\boldsymbol{c}$ 的模 $|\boldsymbol{c}| = |\boldsymbol{a} \times \boldsymbol{b}| = |\boldsymbol{a}||\boldsymbol{b}|\sin\widehat{\boldsymbol{a},\boldsymbol{b}}$，$\boldsymbol{c}$ 的方向同时垂直于 $\boldsymbol{a}$ 和 $\boldsymbol{b}$，且 $\boldsymbol{a},\boldsymbol{b},\boldsymbol{c}$ 符合右手法则. 特别地，向量积 $\boldsymbol{a} \times \boldsymbol{b}$ 的模 $|\boldsymbol{a} \times \boldsymbol{b}| = |\boldsymbol{a}||\boldsymbol{b}|\sin\widehat{\boldsymbol{a},\boldsymbol{b}}$ 在几何上等于以 $\boldsymbol{a},\boldsymbol{b}$ 为邻边的平行四边形的面积.
\end{dl}
\begin{center}
\begin{tikzpicture}[mpdiagram,scale=0.95]
  \coordinate (O) at (0,0);
  \coordinate (A) at (3.0,0.55);
  \coordinate (B) at (1.05,1.65);
  \coordinate (C) at ($(A)+(B)$);
  \path[mpregion] (O) -- (A) -- (C) -- (B) -- cycle;
  \draw[mpvector] (O) -- (A) node[below right,mplabel] {$\boldsymbol a$};
  \draw[mpvector] (O) -- (B) node[above left,mplabel] {$\boldsymbol b$};
  \draw[mpguide] (A) -- (C);
  \draw[mpguide] (B) -- (C);
  \draw[mpaccent] (0.35,0.25) -- (0.35,2.25) node[above,mplabel] {$\boldsymbol a\times\boldsymbol b$};
  \node[mplabel] at (2.25,1.15) {$|\boldsymbol a\times\boldsymbol b|$};
\end{tikzpicture}
\diagramcaption{向量积的方向垂直于两向量所在平面，模等于平行四边形面积。}
\end{center}
\begin{zs}
	    运算律：对任意的向量 $\boldsymbol{a},\boldsymbol{b},\boldsymbol{c}$ 和实数 $\lambda$，
\[
\boldsymbol{a} \times \boldsymbol{b} = -\boldsymbol{b} \times \boldsymbol{a},\quad
\boldsymbol{a} \times (\boldsymbol{b} + \boldsymbol{c}) = \boldsymbol{a} \times \boldsymbol{b} + \boldsymbol{a} \times \boldsymbol{c},\quad
\lambda(\boldsymbol{a} \times \boldsymbol{b}) = (\lambda\boldsymbol{a}) \times \boldsymbol{b} = \boldsymbol{a} \times (\lambda\boldsymbol{b}).
\]
\end{zs}
\begin{zs}
    应用：判定非零向量平行.
\[
\boldsymbol{a} \parallel \boldsymbol{b}
\Leftrightarrow \boldsymbol{a} \times \boldsymbol{b} = \boldsymbol{0}
\Leftrightarrow
\begin{cases}
x_1y_2-y_1x_2=0,\\
x_1z_2-z_1x_2=0,\\
y_1z_2-z_1y_2=0.
\end{cases}
\]
若各对应分量都不为 $0$，也可写成坐标成比例；遇到分量为 $0$ 时，用向量积判定最稳。

\end{zs}

\begin{dl}{ 混合积}{}
设向量 $\boldsymbol{a} = (x_1,y_1,z_1)$，$\boldsymbol{b} = (x_2,y_2,z_2)$，$\boldsymbol{c} = (x_3,y_3,z_3)$，则
\[
[\boldsymbol{a}\ \boldsymbol{b}\ \boldsymbol{c}] = (\boldsymbol{a} \times \boldsymbol{b}) \cdot \boldsymbol{c} =
\begin{vmatrix}
x_1 & y_1 & z_1 \\
x_2 & y_2 & z_2 \\
x_3 & y_3 & z_3
\end{vmatrix}.
\]   
\end{dl}
\begin{zs}
    运算律：
\[
[\boldsymbol{a}\ \boldsymbol{b}\ \boldsymbol{c}] = [\boldsymbol{b}\ \boldsymbol{c}\ \boldsymbol{a}] = [\boldsymbol{c}\ \boldsymbol{a}\ \boldsymbol{b}].
\]
\end{zs}

\begin{zs}
    应用：判定非零向量共面.
若 $[\boldsymbol{a}\ \boldsymbol{b}\ \boldsymbol{c}] = 0$，则 $\boldsymbol{a},\boldsymbol{b},\boldsymbol{c}$ 共面；若 $[\boldsymbol{a}\ \boldsymbol{b}\ \boldsymbol{c}] \neq 0$，则 $\boldsymbol{a},\boldsymbol{b},\boldsymbol{c}$ 不共面，此时，$|[\boldsymbol{a}\ \boldsymbol{b}\ \boldsymbol{c}]|$ 表示以 $\boldsymbol{a},\boldsymbol{b},\boldsymbol{c}$ 为相邻棱的平行六面体的体积.
\end{zs}

\begin{problem}
    设向量 $\boldsymbol{a} = (2,-3,1)$, $\boldsymbol{b} = (1,-2,5)$, 已知向量 $\boldsymbol{c}$ 满足 $\boldsymbol{c} \perp \boldsymbol{a}$, $\boldsymbol{c} \perp \boldsymbol{b}$, 且 $\boldsymbol{c} \cdot (\boldsymbol{i} + 2\boldsymbol{j} - 7\boldsymbol{k}) = 10$, 求 $\boldsymbol{c}$.
\end{problem}
\begin{note}
     利用向量的概念、运算、向量平行及垂直等关系求解相关问题:
对于向量的模、方向余弦（方向角）以及向量的夹角等计算问题，依据定义就可直接计算。利用向量的数量积、向量积以及混合积性质可判断向量是否共线、垂直以及向量组是否共面等空间位置关系.
\end{note}
 \begin{solution}
     因为 $\boldsymbol{c} \perp \boldsymbol{a}$, $\boldsymbol{c} \perp \boldsymbol{b}$, 所以 $\boldsymbol{c}$ 与 $\boldsymbol{a} \times \boldsymbol{b}$ 共线, 即存在实数 $\lambda$, 使得
\[\boldsymbol{c} = \lambda (\boldsymbol{a} \times \boldsymbol{b}) = \lambda
\begin{vmatrix}
\boldsymbol{i} & \boldsymbol{j} & \boldsymbol{k} \\
2 & -3 & 1 \\
1 & -2 & 5
\end{vmatrix}
= -\lambda(13\boldsymbol{i} + 9\boldsymbol{j} + \boldsymbol{k}),\]
又因为\[\boldsymbol{c} \cdot (\boldsymbol{i} + 2\boldsymbol{j} - 7\boldsymbol{k}) = 10,\]
即\[
-\lambda(13\boldsymbol{i} + 9\boldsymbol{j} + \boldsymbol{k}) \cdot (\boldsymbol{i} + 2\boldsymbol{j} - 7\boldsymbol{k}) = 10,\]
计算数量积：\[
-\lambda(13 \times 1 + 9 \times 2 + 1 \times (-7)) = -\lambda(13 + 18 - 7) = -24\lambda = 10,\]
可解得\[
\lambda = -\frac{5}{12}.\]
将 $\lambda = -\frac{5}{12}$ 代入, 可得所求向量\[
\boldsymbol{c} = \frac{5}{12}(13\boldsymbol{i} + 9\boldsymbol{j} + \boldsymbol{k}).\]

 \end{solution}
    \begin{problem}
        求以 $A(1,1,1)$, $B(2,3,4)$ 和 $C(4,3,2)$ 为顶点的三角形的面积;
    \end{problem}
\begin{solution}
    \[
\overrightarrow{AB} = (1,2,3),\quad \overrightarrow{AC} = (3,2,1),
\]
所以
\[
\overrightarrow{AB} \times \overrightarrow{AC} =
\begin{vmatrix}
\boldsymbol{i} & \boldsymbol{j} & \boldsymbol{k} \\
1 & 2 & 3 \\
3 & 2 & 1
\end{vmatrix}
= -4\boldsymbol{i} + 8\boldsymbol{j} - 4\boldsymbol{k}.
\]
由向量积模的几何意义, 可知所求 $\triangle ABC$ 的面积
\[
S = \frac{1}{2} |\overrightarrow{AB} \times \overrightarrow{AC}|
= \frac{1}{2} \sqrt{(-4)^2 + 8^2 + (-4)^2}
= 2\sqrt{6}.
\]
\end{solution}
\begin{problem}
    设向量 $\boldsymbol{a},\boldsymbol{b},\boldsymbol{c}$ 共面，$\boldsymbol{a} = (0,1,2)$, $\boldsymbol{b} = (-2,1,3)$, $\boldsymbol{c}$ 在 $\boldsymbol{a}$ 上的投影为 $\sqrt{5}$，$\boldsymbol{c}$ 在 $\boldsymbol{b}$ 上的投影为 $-\sqrt{14}$，求 $\boldsymbol{c}$.
\end{problem}
\begin{note}
    根据投影的计算公式以及向量共面的充要条件求解.
\end{note}
 \begin{solution}
     设 $\boldsymbol{c} = (x,y,z)$，由于 $\boldsymbol{c}$ 在 $\boldsymbol{a}$ 上的投影为 $\sqrt{5}$，所以
\[
\text{Prj}_{\boldsymbol{a}} \boldsymbol{c} = \frac{\boldsymbol{a} \cdot \boldsymbol{c}}{|\boldsymbol{a}|} = \frac{y + 2z}{\sqrt{5}} = \sqrt{5},
\]
即
\[
y + 2z = 5.
\]

同理，由 $\boldsymbol{c}$ 在 $\boldsymbol{b}$ 上的投影为 $-\sqrt{14}$，有
\[
\text{Prj}_{\boldsymbol{b}} \boldsymbol{c} = \frac{\boldsymbol{b} \cdot \boldsymbol{c}}{|\boldsymbol{b}|} = \frac{-2x + y + 3z}{\sqrt{14}} = -\sqrt{14},
\]
即
\[
-2x + y + 3z = -14.
\]

又因为向量 $\boldsymbol{a},\boldsymbol{b},\boldsymbol{c}$ 共面，所以 $(\boldsymbol{a} \times \boldsymbol{b}) \cdot \boldsymbol{c} = 0$，即
\[
\begin{vmatrix}
0 & 1 & 2 \\
-2 & 1 & 3 \\
x & y & z
\end{vmatrix} = 0,
\]
按第一行展开：
\[
0 \cdot \begin{vmatrix}1 & 3 \\ y & z\end{vmatrix}
- 1 \cdot \begin{vmatrix}-2 & 3 \\ x & z\end{vmatrix}
+ 2 \cdot \begin{vmatrix}-2 & 1 \\ x & y\end{vmatrix} = 0,
\]
化简得：
\[
-(-2z - 3x) + 2(-2y - x) = 0 \implies 3x + 2z - 4y - 2x = 0 \implies x - 4y + 2z = 0.
\]

联立上述三个方程：
\[
\begin{cases}
y + 2z = 5, \\
-2x + y + 3z = -14, \\
x - 4y + 2z = 0,
\end{cases}
\]
解得 $x = 10$, $y = 3$, $z = 1$，故所求向量 $\boldsymbol{c} = (10,3,1)$.
 \end{solution}
\subsection{第二节 平面与直线}
\subsubsection{平面方程}
\begin{zs}
    点法式: 设平面过点 $M_0(x_0,y_0,z_0)$，法向量为 $\boldsymbol{n}=(A,B,C)$，则平面方程为  \[  A(x-x_0)+B(y-y_0)+C(z-z_0)=0.  \]
\end{zs}

\begin{zs}
    三点式: 设平面过不共线三点 $M_i(x_i,y_i,z_i)\ (i=1,2,3)$，则平面方程为
  \[
  \begin{vmatrix}
  x-x_1 & y-y_1 & z-z_1 \\
  x_2-x_1 & y_2-y_1 & z_2-z_1 \\
  x_3-x_1 & y_3-y_1 & z_3-z_1
  \end{vmatrix}=0,
  \]
  其中法向量为 $\boldsymbol{n} = \overrightarrow{M_1M_2} \times \overrightarrow{M_1M_3}$.
\end{zs}
\begin{zs}
    一般式: $Ax+By+Cz+D=0$，其中 $\boldsymbol{n}=(A,B,C)$ 为平面的法向量.
\end{zs}
\begin{zs}
    截距式: $\dfrac{x}{a}+\dfrac{y}{b}+\dfrac{z}{c}=1$，其中 $a,b,c$ 分别为平面在 $x$ 轴、$y$ 轴、$z$ 轴上的截距.
\end{zs}
\begin{zs}
    特殊平面:

  过原点的平面：$Ax+By+Cz=0$
  
  平行于 $x$ 轴（$y$ 轴、$z$ 轴）：$By+Cz+D=0$（$Ax+Cz+D=0$、$Ax+By+D=0$）
  
平行于 $xOy$ 面（$yOz$ 面、$zOx$ 面）：$Cz+D=0$（$Ax+D=0$、$By+D=0$）


 三个坐标面：$xOy$ 面 $z=0$，$yOz$ 面 $x=0$，$zOx$ 面 $y=0$

\end{zs}
\subsubsection{平面的位置关系}
 设平面 $\Pi_1: A_1x+B_1y+C_1z+D_1=0$，$\Pi_2: A_2x+B_2y+C_2z+D_2=0$，则
\begin{zs}
   平行或重合: $\Pi_1$ 与 $\Pi_2$ 的法向量平行，即 $\boldsymbol{n}_1\times \boldsymbol{n}_2=\boldsymbol{0}$。
   若进一步有 $(A_1,B_1,C_1,D_1)$ 与 $(A_2,B_2,C_2,D_2)$ 成比例，则两平面重合；否则为平行不重合。
\end{zs}
\begin{zs}
    垂直: $\Pi_1 \perp \Pi_2 \Leftrightarrow A_1A_2+B_1B_2+C_1C_2=0$
\end{zs}
\begin{zs}
    夹角: $\cos\theta=\dfrac{|\boldsymbol{n}_1\cdot\boldsymbol{n}_2|}{|\boldsymbol{n}_1||\boldsymbol{n}_2|}=\dfrac{|A_1A_2+B_1B_2+C_1C_2|}{\sqrt{A_1^2+B_1^2+C_1^2}\sqrt{A_2^2+B_2^2+C_2^2}}$
\end{zs}
\subsubsection{直线方程}
直线过点 $M_0(x_0,y_0,z_0)$，方向向量为 $\boldsymbol{s}=(m,n,p)$ 的直线方程为
\begin{zs}
    点向式（对称式）: 
  \[
  \dfrac{x-x_0}{m}=\dfrac{y-y_0}{n}=\dfrac{z-z_0}{p}.
  \]
\end{zs}
\begin{zs}
    参数式:
  \[
  \begin{cases}
  x=x_0+mt, \\
  y=y_0+nt, \\
  z=z_0+pt,
  \end{cases}
  \quad t\in\mathbb{R}.
  \]
\end{zs}
\begin{zs}
    一般式（两平面的交线）:
  \[
  \begin{cases}
  A_1x+B_1y+C_1z+D_1=0, \\
  A_2x+B_2y+C_2z+D_2=0,
  \end{cases}
  \]
  方向向量为 $\boldsymbol{s} = \boldsymbol{n}_1 \times \boldsymbol{n}_2$，其中 $\boldsymbol{n}_1=(A_1,B_1,C_1),\ \boldsymbol{n}_2=(A_2,B_2,C_2)$ 为两平面的法向量.

注意：直线的一般式方程不唯一
\end{zs}
\begin{zs}
    两点式: 过点 $M_1(x_1,y_1,z_1)$ 和 $M_2(x_2,y_2,z_2)$ 的直线方程为
  \[
  \dfrac{x-x_1}{x_2-x_1}=\dfrac{y-y_1}{y_2-y_1}=\dfrac{z-z_1}{z_2-z_1},
  \]
  方向向量为 $\boldsymbol{s} = \overrightarrow{M_1M_2}$.
\end{zs}
\subsubsection{直线的位置关系}
设直线 $L_1: \dfrac{x-x_1}{m_1}=\dfrac{y-y_1}{n_1}=\dfrac{z-z_1}{p_1}$，$L_2: \dfrac{x-x_2}{m_2}=\dfrac{y-y_2}{n_2}=\dfrac{z-z_2}{p_2}$，则
\begin{zs}
    平行: $L_1 \parallel L_2 \Leftrightarrow \boldsymbol{s}_1\times\boldsymbol{s}_2=\boldsymbol{0}$，即两条直线的方向向量平行。
\end{zs}
\begin{zs}
    垂直: $L_1 \perp L_2 \Leftrightarrow m_1m_2+n_1n_2+p_1p_2=0$
\end{zs}
\begin{zs}
    夹角: $\cos\theta=\dfrac{|\boldsymbol{s}_1\cdot\boldsymbol{s}_2|}{|\boldsymbol{s}_1||\boldsymbol{s}_2|}=\dfrac{|m_1m_2+n_1n_2+p_1p_2|}{\sqrt{m_1^2+n_1^2+p_1^2}\sqrt{m_2^2+n_2^2+p_2^2}}$
\end{zs}
\subsubsection{直线与平面的位置关系}
设直线 $L: \dfrac{x-x_0}{m}=\dfrac{y-y_0}{n}=\dfrac{z-z_0}{p}$，平面 $\Pi: Ax+By+Cz+D=0$，则
\begin{zs}
    平行: $L \parallel \Pi \Leftrightarrow Am+Bn+Cp=0$
\end{zs}
\begin{zs}
    垂直: $L \perp \Pi \Leftrightarrow \boldsymbol{s}\parallel\boldsymbol{n}$，即 $\boldsymbol{s}\times\boldsymbol{n}=\boldsymbol{0}$。
\end{zs}
\begin{zs}
    夹角: $\sin\varphi=\dfrac{|\boldsymbol{s}\cdot\boldsymbol{n}|}{|\boldsymbol{s}||\boldsymbol{n}|}=\dfrac{|Am+Bn+Cp|}{\sqrt{A^2+B^2+C^2}\sqrt{m^2+n^2+p^2}}$
\end{zs}
\subsubsection{距离}
\begin{center}
\begin{tikzpicture}[mpdiagram,scale=0.95]
  \coordinate (A) at (-2.6,-0.55);
  \coordinate (B) at (2.4,-0.20);
  \coordinate (C) at (1.55,1.05);
  \coordinate (D) at (-3.05,0.65);
  \coordinate (H) at (-0.45,0.24);
  \coordinate (P) at (-0.45,2.08);
  \path[mpsoftregion] (A) -- (B) -- (C) -- (D) -- cycle;
  \node[mplabel] at (2.05,0.72) {$\Pi$};
  \fill[MaterialRed] (P) circle (1.25pt) node[above,mplabel] {$M_0$};
  \fill[MaterialBlueGrey] (H) circle (1pt) node[below,mplabel] {$H$};
  \draw[mpguide] (P) -- (H);
  \draw[mpaccent] (H) -- ++(0,1.28) node[right,mplabel] {$\boldsymbol n$};
  \draw[decorate,decoration={brace,amplitude=4pt},MaterialRed]
    ($(H)+(-0.18,0)$) -- ($(P)+(-0.18,0)$)
    node[midway,left=5pt,mplabel] {$d$};
  \draw[MaterialBlueGrey!65] ($(H)+(0.20,0)$) -- ++(0,0.20) -- ++(-0.20,0);
\end{tikzpicture}
\diagramcaption{点到平面的距离，是点沿平面法向量方向到平面的最短距离。}
\end{center}
\begin{zs}
点到平面距离：
    点 $M_0(x_0,y_0,z_0)$ 到平面 $Ax+By+Cz+D=0$ 的距离为
\[
d=\dfrac{|Ax_0+By_0+Cz_0+D|}{\sqrt{A^2+B^2+C^2}}.
\]
\end{zs}
\begin{zs}
    点到直线距离：点 $M_0(x_0,y_0,z_0)$ 到直线 $L: \dfrac{x-x_1}{m}=\dfrac{y-y_1}{n}=\dfrac{z-z_1}{p}$ 的距离为
\[
d=\dfrac{|\overrightarrow{M_0M_1} \times \boldsymbol{s}|}{|\boldsymbol{s}|},
\]
其中 $M_1(x_1,y_1,z_1) \in L$，$\boldsymbol{s}=(m,n,p)$ 为直线 $L$ 的方向向量.

\end{zs}
\begin{zs}
    两异面直线距离：设异面直线 $L_1: \dfrac{x-x_1}{m_1}=\dfrac{y-y_1}{n_1}=\dfrac{z-z_1}{p_1}$ 与 $L_2: \dfrac{x-x_2}{m_2}=\dfrac{y-y_2}{n_2}=\dfrac{z-z_2}{p_2}$，距离为
\[
d=\dfrac{|\overrightarrow{M_1M_2} \cdot (\boldsymbol{s}_1 \times \boldsymbol{s}_2)|}{|\boldsymbol{s}_1 \times \boldsymbol{s}_2|},
\]
其中 $M_1(x_1,y_1,z_1)\in L_1,\ M_2(x_2,y_2,z_2)\in L_2$，方向向量为 $\boldsymbol{s}_1=(m_1,n_1,p_1),\ \boldsymbol{s}_2=(m_2,n_2,p_2)$.
\end{zs}
\subsubsection{平面束}
\begin{zs}
    设直线 $L$ 的方程为
\[
\begin{cases}
A_1x+B_1y+C_1z+D_1=0, \\
A_2x+B_2y+C_2z+D_2=0,
\end{cases}
\]
则过直线 $L$（除平面 $A_1x+B_1y+C_1z+D_1=0$ 外）的平面束的一般表达式为
\[
\lambda(A_1x+B_1y+C_1z+D_1)+(A_2x+B_2y+C_2z+D_2)=0,
\]
\end{zs}
\begin{problem}
   求过直线 $L: \begin{cases} 4x+2y+3z-6=0, \\ 2x+y=0 \end{cases}$ 且与球面 $x^2+y^2+z^2=4$ 相切的平面 $\Pi$ 的方程. 
\end{problem}
\begin{note}
    要计算平面的方程，一般情况下，需知道平面上一点的坐标和平面的法向量，以此利用点法式获得平面的方程。对于一些特殊情况，也可用平面束方法来计算.
\end{note}
\begin{solution}
    利用平面束法建立平面方程。由于所求平面 $\Pi$ 过直线 $L$，所以可设所求平面方程为
\[
4x + 2y + 3z - 6 + \lambda(2x + y) = 0,
\]
即
\[
(4+2\lambda)x + (2+\lambda)y + 3z - 6 = 0.
\]
其法向量 $\boldsymbol{n}=(4+2\lambda,2+\lambda,3)$。由题意，球心 $(0,0,0)$ 到此平面的距离为 $2$，即
\[
\dfrac{|(4+2\lambda)\cdot0 + (2+\lambda)\cdot0 + 3\cdot0 - 6|}{\sqrt{(4+2\lambda)^2+(2+\lambda)^2+3^2}} = 2,
\]
解得 $\lambda = -2$.将其代入，可得所求平面 $\Pi$ 的方程为 $z=2$.
\end{solution}
\begin{problem}
    用对称式及参数式表示直线：\[
\begin{cases}
x + y + z + 1 = 0, \\
2x - y + 3z + 4 = 0.
\end{cases}
\]
\end{problem}
\begin{solution}
    先在直线上任取一点，不妨设 $x=1$，解方程组
\[
\begin{cases}
y + z = -2, \\
y - 3z = 6,
\end{cases}
\]
得 $y=0, z=-2$，即 $M_0(1,0,-2)$ 为直线上一点。再计算直线的方向向量，可取直线的方向向量
\[
\boldsymbol{s} = \begin{vmatrix}
\boldsymbol{i} & \boldsymbol{j} & \boldsymbol{k} \\
1 & 1 & 1 \\
2 & -1 & 3
\end{vmatrix} = 4\boldsymbol{i} - \boldsymbol{j} - 3\boldsymbol{k}.
\]
故直线的对称式方程为
\[
\dfrac{x-1}{4} = \dfrac{y}{-1} = \dfrac{z+2}{-3}.
\]
若令 $\dfrac{x-1}{4} = \dfrac{y}{-1} = \dfrac{z+2}{-3} = t$，可得直线的参数式方程为
\[
\begin{cases}
x = 1 + 4t, \\
y = -t, \\
z = -2 - 3t,
\end{cases}
\quad (-\infty < t < +\infty).
\]
\end{solution}
\subsection{第三节 曲面与曲线}
\subsubsection{曲面方程}
\begin{center}
\begin{tikzpicture}[mpdiagram,scale=0.88]
  \begin{scope}[xshift=-3.2cm]
    \draw[mpboundary] (0,0) ellipse (1.10 and 0.36);
    \draw[mpboundary] (0,2.25) ellipse (1.10 and 0.36);
    \draw[mpboundary] (-1.10,0) -- (-1.10,2.25);
    \draw[mpboundary] (1.10,0) -- (1.10,2.25);
    \draw[mpguide] (-0.55,0.12) -- (-0.55,2.37);
    \draw[mpguide] (0.55,-0.12) -- (0.55,2.13);
    \node[mplabel] at (0,-0.78) {柱面};
    \node[mplabel] at (0,2.86) {母线平行};
  \end{scope}
  \begin{scope}[xshift=3.1cm]
    \draw[mpaxis] (0,-0.15) -- (0,2.85) node[above,mplabel] {旋转轴};
    \draw[mpguide] (-1.55,0.28) .. controls (-0.65,0.42) and (-0.25,1.12) .. (-0.85,2.30);
    \draw[mpguide] (1.55,0.28) .. controls (0.65,0.42) and (0.25,1.12) .. (0.85,2.30);
    \draw[mpboundary] (0,0.28) ellipse (1.55 and 0.28);
    \draw[mpboundary] (0,1.25) ellipse (0.88 and 0.20);
    \draw[mpboundary] (0,2.30) ellipse (0.85 and 0.18);
    \draw[mpaccent] (0.35,0.55) arc[start angle=-105,end angle=205,x radius=0.62,y radius=0.25];
    \node[mplabel] at (0,-0.78) {旋转曲面};
    \node[mplabel] at (1.75,1.95) {母线旋转};
  \end{scope}
\end{tikzpicture}
\diagramcaption{柱面看母线是否平行；旋转曲面看曲线到旋转轴的距离如何替换。}
\end{center}
\begin{dy}{柱面}{}
平行于定直线 $L$ 并沿给定曲线 $C$ 移动的直线所形成的曲面称为柱面，定曲线 $C$ 叫做柱面的准线，动直线叫做柱面的母线。

在空间直角坐标系中，缺少某一个变量的方程表示母线平行于某坐标轴的柱面。例如，$F(x,y)=0$ 表示准线为
\[
\begin{cases}
F(x,y)=0, \\
z=0
\end{cases}
\]
且母线平行于 $z$ 轴的柱面.
\end{dy}
\begin{zs}
    常见的柱面有：
 圆柱面：$x^2+y^2=R^2$
 
抛物柱面：$y^2=2px$

双曲柱面：$\dfrac{x^2}{a^2}-\dfrac{y^2}{b^2}=1$

 椭圆柱面：$\dfrac{x^2}{a^2}+\dfrac{y^2}{b^2}=1$
\end{zs}
\begin{dy}{旋转曲面}{}
曲线 $C$ 绕定直线 $L$ 旋转所形成的曲面称为旋转曲面，曲线 $C$ 叫做旋转曲面的母线，定直线 $L$ 叫做旋转曲面的轴（或旋转轴）。

例如，$xOy$ 坐标面上的曲线
\[
\begin{cases}
f(x,y)=0, \\
z=0
\end{cases}
\]
绕 $x$ 轴旋转所形成的旋转曲面的方程为 $f(x,\pm\sqrt{y^2+z^2})=0$.
\end{dy}
    \begin{zs}
        圆锥面：$z^2=x^2+y^2$，旋转轴为 $z$ 轴

 旋转抛物面：$z=x^2+y^2$，旋转轴为 $z$ 轴

单叶旋转双曲面：$\dfrac{x^2+z^2}{a^2}-\dfrac{y^2}{b^2}=1$，旋转轴为 $y$ 轴

双叶旋转双曲面：$\dfrac{x^2}{a^2}-\dfrac{y^2+z^2}{b^2}=-1$，旋转轴为 $x$ 轴
    \end{zs}
\subsubsection*{ 几种特殊的二次曲面}
\begin{zs}
    椭球面：$\dfrac{x^2}{a^2}+\dfrac{y^2}{b^2}+\dfrac{z^2}{c^2}=1$.

当 $a,b,c$ 中有两个相等时，如 $\dfrac{x^2}{a^2}+\dfrac{y^2}{a^2}+\dfrac{z^2}{c^2}=1$，均表示旋转椭球面.

特别地，当 $a=b=c=R$ 时，方程 $x^2+y^2+z^2=R^2$ 表示球面.
\end{zs}
\begin{zs}
    抛物面：分为椭圆抛物面和双曲抛物面.

椭圆抛物面：$\dfrac{x^2}{a^2}+\dfrac{y^2}{b^2}=\pm z$，

当 $a=b$ 时，$\dfrac{x^2}{a^2}+\dfrac{y^2}{a^2}=\pm z$ 表示旋转抛物面

双曲抛物面(马鞍面)：$\dfrac{x^2}{a^2}-\dfrac{y^2}{b^2}=\pm z$
\end{zs}
\begin{zs}
    双曲面：分为单叶双曲面和双叶双曲面.
单叶双曲面：$\dfrac{x^2}{a^2}+\dfrac{y^2}{b^2}-\dfrac{z^2}{c^2}=1$，

当 $a=b$ 时，$\dfrac{x^2}{a^2}+\dfrac{y^2}{a^2}-\dfrac{z^2}{c^2}=1$ 表示单叶旋转双曲面

双叶双曲面：$\dfrac{x^2}{a^2}+\dfrac{y^2}{b^2}-\dfrac{z^2}{c^2}=-1$，

当 $a=b$ 时，$\dfrac{x^2}{a^2}+\dfrac{y^2}{a^2}-\dfrac{z^2}{c^2}=-1$ 表示双叶旋转双曲面
\end{zs}
\begin{zs}
    椭圆锥面：$\dfrac{x^2}{a^2}+\dfrac{y^2}{b^2}-\dfrac{z^2}{c^2}=0$. 

当 $a=b$ 时，$\dfrac{x^2}{a^2}+\dfrac{y^2}{a^2}-\dfrac{z^2}{c^2}=0$ 表示圆锥面.
\end{zs}
\subsubsection{曲线方程}
\begin{zs}
    一般式方程：
  \[
  \begin{cases}
  F(x,y,z)=0, \\
  G(x,y,z)=0.
  \end{cases}
  \]
  注：空间曲线的一般式方程不唯一.
\end{zs}
\begin{zs}
    曲线的参数式方程：
  \[
  \begin{cases}
  x=x(t), \\
  y=y(t), \\
  z=z(t),
  \end{cases}
  \quad (t\in I\subseteq\mathbb{R}).
  \]
  例如，圆柱螺线
  \[
  \begin{cases}
  x=a\cos\omega t, \\
  y=a\sin\omega t, \\
  z=bt
  \end{cases}
  \quad (t\in\mathbb{R});
  \]
\end{zs}
\begin{zs}
    空间曲线在坐标面上的投影：设曲线 $C$:
\[
\begin{cases}
F(x,y,z)=0, \\
G(x,y,z)=0,
\end{cases}
\]
在 $xOy$ 面上的投影：消去方程组中的变量 $z$，得投影柱面 $H(x,y)=0$，则曲线 $C$ 在 $xOy$ 面上的投影曲线方程为
\[
\begin{cases}
H(x,y)=0, \\
z=0.
\end{cases}
\]
类似可得，曲线 $C$ 在 $yOz$ 面、$zOx$ 面上的投影.
\end{zs}
\begin{problem}
    求准线为 $C: \begin{cases} x^2+y^2-z^2=1, \\ x+y+z=1 \end{cases}$，母线平行于 $z$ 轴的柱面方程;
\end{problem}
\begin{note}
    根据柱面定义或者旋转曲面形成规律写出曲面方程即可.例如，准线为 $xOy$ 面上的曲线 $C: \begin{cases} f(x,y)=0, \\ z=0 \end{cases}$，母线平行于 $z$ 轴的柱面方程为 $f(x,y)=0$；曲线 $C: \begin{cases} f(y,z)=0, \\ x=0 \end{cases}$ 绕 $y$ 轴旋转形成的旋转曲面的方程为 $f(y,\pm\sqrt{x^2+z^2})=0$，其他情形类似可得.
\end{note}
\begin{solution}
    根据母线平行于 $z$ 轴的特点，从 $C: \begin{cases} x^2+y^2-z^2=1, \\ x+y+z=1 \end{cases}$ 中消去 $z$，得柱面方程为 $x^2+y^2-(1-x-y)^2=1$.
\end{solution}
\begin{problem}
    椭球面 $S_1$ 是由椭圆 $\dfrac{x^2}{4}+\dfrac{y^2}{3}=1$ 绕 $x$ 轴旋转而成，圆锥面 $S_2$ 是过点 $(4,0)$ 且与椭圆 $\dfrac{x^2}{4}+\dfrac{y^2}{3}=1$ 相切的直线绕 $x$ 轴旋转而成.

(1) 求椭球面 $S_1$ 和圆锥面 $S_2$ 的方程;

(2) 求椭球面 $S_1$ 和圆锥面 $S_2$ 围成的立体的体积.

\end{problem}
\begin{note}
    旋转曲面 $S_1$ 与 $S_2$ 均由坐标面上的曲线绕坐标轴旋转形成，所以根据旋转曲面形成规律可直接写出旋转曲面的方程.所围立体体积可看成两个旋转体体积之差.
\end{note}
\begin{solution}
    (1) 旋转椭球面 $S_1$ 的方程为
\[
\dfrac{x^2}{4}+\dfrac{y^2+z^2}{3}=1.
\]

下面求圆锥面 $S_2$ 的方程.首先，求过点 $(4,0)$ 且与椭圆 $\dfrac{x^2}{4}+\dfrac{y^2}{3}=1$ 相切的直线的方程。设此椭圆上的一点为 $M_0(x_0,y_0)$，则点 $M_0(x_0,y_0)$ 处切线的方程为
\[
\dfrac{xx_0}{4}+\dfrac{yy_0}{3}=1.
\]
要求该切线过点 $(4,0)$ 且 $M_0(x_0,y_0)$ 为椭圆 $\dfrac{x^2}{4}+\dfrac{y^2}{3}=1$ 上的点，即满足
\[
\begin{cases}
\dfrac{4x_0}{4}+\dfrac{0\cdot y_0}{3}=1, \\
\dfrac{x_0^2}{4}+\dfrac{y_0^2}{3}=1.
\end{cases}
\]
解得切点为 $\left(1,\pm\dfrac{3}{2}\right)$，斜率为 $k=\mp\dfrac{1}{2}$，即切线方程为 $y=-\dfrac{1}{2}x+2$ 或 $y=\dfrac{1}{2}x-2$.

其次，求圆锥面 $S_2$ 的方程.圆锥面 $S_2$ 可看成由直线 $y=-\dfrac{1}{2}x+2$ 或 $y=\dfrac{1}{2}x-2$ 绕 $x$ 轴旋转而成，其方程为
\[
\pm\sqrt{y^2+z^2}=-\dfrac{1}{2}x+2,
\]
即
\[
y^2+z^2=\left(\dfrac{1}{2}x-2\right)^2.
\]

(2) 设 $V_1$ 表示由直线段 $y_1=-\dfrac{1}{2}x+2\ (1\leqslant x\leqslant4)$, $y=0$ 及 $x=1$ 所围三角形绕 $x$ 轴旋转形成的旋转体体积，$V_2$ 表示由曲线段 $y_2=\sqrt{3\left(1-\dfrac{x^2}{4}\right)}\ (1\leqslant x\leqslant2)$, $y=0$ 及 $x=1$ 所围曲边三角形绕 $x$ 轴旋转形成的旋转体体积，则所求立体体积为 $V_1$ 与 $V_2$ 的差.

\[
V_1 = \int_{1}^{4} \pi y_1^2 \mathrm{d}x = \int_{1}^{4} \pi \left(-\dfrac{1}{2}x+2\right)^2 \mathrm{d}x = \pi\left(\dfrac{1}{12}x^3 - x^2 + 4x\right)\bigg|_{1}^{4} = \dfrac{9}{4}\pi,
\]
\[
V_2 = \int_{1}^{2} \pi y_2^2 \mathrm{d}x = \int_{1}^{2} \pi \left(3-\dfrac{3}{4}x^2\right) \mathrm{d}x = \pi\left(3x-\dfrac{1}{4}x^3\right)\bigg|_{1}^{2} = \dfrac{5}{4}\pi.
\]

所以椭球面 $S_1$ 和圆锥面 $S_2$ 围成的立体的体积为
\[
V = V_1 - V_2 = \dfrac{9}{4}\pi - \dfrac{5}{4}\pi = \pi.
\]

\paragraph{注}
$V_1$ 也可由圆锥体积公式求得.
\end{solution}

\begin{warningbox}{第五章易错点}
\begin{itemize}
  \item 判断向量平行时，优先用向量积为零；只有对应分量都不为零时，才直接写成分量比相等。
  \item 写平面方程先找法向量，写直线方程先找方向向量；不要把二者混用。
  \item 对称式中若某个方向分量为 $0$，表示对应坐标保持不变，不表示分母可以随意约去。
  \item 柱面方程常缺少一个变量；旋转曲面要把到旋转轴的距离替换正确。
\end{itemize}
\end{warningbox}
